% Options for packages loaded elsewhere
% Options for packages loaded elsewhere
\PassOptionsToPackage{unicode}{hyperref}
\PassOptionsToPackage{hyphens}{url}
\PassOptionsToPackage{dvipsnames,svgnames,x11names}{xcolor}
%
\documentclass[
  letterpaper,
  DIV=11,
  numbers=noendperiod]{scrartcl}
\usepackage{xcolor}
\usepackage{amsmath,amssymb}
\setcounter{secnumdepth}{-\maxdimen} % remove section numbering
\usepackage{iftex}
\ifPDFTeX
  \usepackage[T1]{fontenc}
  \usepackage[utf8]{inputenc}
  \usepackage{textcomp} % provide euro and other symbols
\else % if luatex or xetex
  \usepackage{unicode-math} % this also loads fontspec
  \defaultfontfeatures{Scale=MatchLowercase}
  \defaultfontfeatures[\rmfamily]{Ligatures=TeX,Scale=1}
\fi
\usepackage{lmodern}
\ifPDFTeX\else
  % xetex/luatex font selection
\fi
% Use upquote if available, for straight quotes in verbatim environments
\IfFileExists{upquote.sty}{\usepackage{upquote}}{}
\IfFileExists{microtype.sty}{% use microtype if available
  \usepackage[]{microtype}
  \UseMicrotypeSet[protrusion]{basicmath} % disable protrusion for tt fonts
}{}
\makeatletter
\@ifundefined{KOMAClassName}{% if non-KOMA class
  \IfFileExists{parskip.sty}{%
    \usepackage{parskip}
  }{% else
    \setlength{\parindent}{0pt}
    \setlength{\parskip}{6pt plus 2pt minus 1pt}}
}{% if KOMA class
  \KOMAoptions{parskip=half}}
\makeatother
% Make \paragraph and \subparagraph free-standing
\makeatletter
\ifx\paragraph\undefined\else
  \let\oldparagraph\paragraph
  \renewcommand{\paragraph}{
    \@ifstar
      \xxxParagraphStar
      \xxxParagraphNoStar
  }
  \newcommand{\xxxParagraphStar}[1]{\oldparagraph*{#1}\mbox{}}
  \newcommand{\xxxParagraphNoStar}[1]{\oldparagraph{#1}\mbox{}}
\fi
\ifx\subparagraph\undefined\else
  \let\oldsubparagraph\subparagraph
  \renewcommand{\subparagraph}{
    \@ifstar
      \xxxSubParagraphStar
      \xxxSubParagraphNoStar
  }
  \newcommand{\xxxSubParagraphStar}[1]{\oldsubparagraph*{#1}\mbox{}}
  \newcommand{\xxxSubParagraphNoStar}[1]{\oldsubparagraph{#1}\mbox{}}
\fi
\makeatother


\usepackage{longtable,booktabs,array}
\usepackage{calc} % for calculating minipage widths
% Correct order of tables after \paragraph or \subparagraph
\usepackage{etoolbox}
\makeatletter
\patchcmd\longtable{\par}{\if@noskipsec\mbox{}\fi\par}{}{}
\makeatother
% Allow footnotes in longtable head/foot
\IfFileExists{footnotehyper.sty}{\usepackage{footnotehyper}}{\usepackage{footnote}}
\makesavenoteenv{longtable}
\usepackage{graphicx}
\makeatletter
\newsavebox\pandoc@box
\newcommand*\pandocbounded[1]{% scales image to fit in text height/width
  \sbox\pandoc@box{#1}%
  \Gscale@div\@tempa{\textheight}{\dimexpr\ht\pandoc@box+\dp\pandoc@box\relax}%
  \Gscale@div\@tempb{\linewidth}{\wd\pandoc@box}%
  \ifdim\@tempb\p@<\@tempa\p@\let\@tempa\@tempb\fi% select the smaller of both
  \ifdim\@tempa\p@<\p@\scalebox{\@tempa}{\usebox\pandoc@box}%
  \else\usebox{\pandoc@box}%
  \fi%
}
% Set default figure placement to htbp
\def\fps@figure{htbp}
\makeatother





\setlength{\emergencystretch}{3em} % prevent overfull lines

\providecommand{\tightlist}{%
  \setlength{\itemsep}{0pt}\setlength{\parskip}{0pt}}



 


\KOMAoption{captions}{tableheading}
\makeatletter
\@ifpackageloaded{caption}{}{\usepackage{caption}}
\AtBeginDocument{%
\ifdefined\contentsname
  \renewcommand*\contentsname{Table of contents}
\else
  \newcommand\contentsname{Table of contents}
\fi
\ifdefined\listfigurename
  \renewcommand*\listfigurename{List of Figures}
\else
  \newcommand\listfigurename{List of Figures}
\fi
\ifdefined\listtablename
  \renewcommand*\listtablename{List of Tables}
\else
  \newcommand\listtablename{List of Tables}
\fi
\ifdefined\figurename
  \renewcommand*\figurename{Figure}
\else
  \newcommand\figurename{Figure}
\fi
\ifdefined\tablename
  \renewcommand*\tablename{Table}
\else
  \newcommand\tablename{Table}
\fi
}
\@ifpackageloaded{float}{}{\usepackage{float}}
\floatstyle{ruled}
\@ifundefined{c@chapter}{\newfloat{codelisting}{h}{lop}}{\newfloat{codelisting}{h}{lop}[chapter]}
\floatname{codelisting}{Listing}
\newcommand*\listoflistings{\listof{codelisting}{List of Listings}}
\makeatother
\makeatletter
\makeatother
\makeatletter
\@ifpackageloaded{caption}{}{\usepackage{caption}}
\@ifpackageloaded{subcaption}{}{\usepackage{subcaption}}
\makeatother
\usepackage{bookmark}
\IfFileExists{xurl.sty}{\usepackage{xurl}}{} % add URL line breaks if available
\urlstyle{same}
\hypersetup{
  pdftitle={Del monopolio al oligopolio: la liberalización del sistema ferroviario español},
  colorlinks=true,
  linkcolor={blue},
  filecolor={Maroon},
  citecolor={Blue},
  urlcolor={Blue},
  pdfcreator={LaTeX via pandoc}}


\title{Del monopolio al oligopolio: la liberalización del sistema
ferroviario español}
\author{Javier Ribal del Río}
\date{2026-03-29}
\begin{document}
\maketitle


\subsection{Introducción}\label{introducciuxf3n}

\subsubsection{Antecedentes}\label{antecedentes}

En 1941 se crea \emph{Red Nacional de los Ferrocarriles Españoles
(RENFE)} una compañía estatal fundada con el objetivo de estandarizar la
red ferroviaria española actuando como operadora de transporte
ferroviario y como gestora de la infraestructura (vías, estaciones,
red). Esta configuración se corresponde con un monopolio natural,
caracterizado por la existencia de \emph{RENFE} como compañía
monopolística que concentra toda la oferta ferroviaria. Este monopolio
terminaría con el proceso de liberalización del sistema ferroviario
iniciado en la directiva 91/440 de la entonces CEE, en primer lugar
provocó la escisión en 2005 de \emph{RENFE} en \emph{Administrador de
Infraestructuras Ferroviarias (Adif)} (gestión de infraestructuras) y
\emph{RENFE-Operadora} (operadora de la red), y finalmente, en 2020 la
apertura de la red a compañías privadas. La entrada de la operadora
Francesa \emph{OUIGO} en 2021 e \emph{iryo} en 2022, supuso el paso de
un monopolio a un oligopolio.

\subsubsection{Objetivos}\label{objetivos}

El objetivo del presente trabajo académico es analizar la estructura de
mercado del transporte ferroviario español tras su liberalización,
particularmente en el sector de la alta velocidad, evaluando si se
corresponde con un oligopolio. Interpretandolo con modelos del
oligopolio, analizando el grado de concentración de mercado y el
comportamiento microeconómico de las empresas.

\subsubsection{Justificación}\label{justificaciuxf3n}

En la actualidad España mantiene la 2ª red de alta velocidad más extensa
del mundo, por lo que resulta interesante en un contexto con un alto
nivel de utilización de la red, analizar la transición de un monopolio a
un oligopolio y su impacto en la competencia.

\begin{center}\rule{0.5\linewidth}{0.5pt}\end{center}

\subsection{Trabajo restante}\label{trabajo-restante}

\begin{itemize}
\item
  \begin{enumerate}
  \def\labelenumi{\arabic{enumi}.}
  \setcounter{enumi}{1}
  \tightlist
  \item
    DESARROLLO
  \end{enumerate}

  \begin{itemize}
  \tightlist
  \item
    2.1 Estructura del mercado y barreras de entrada:(operadores
    principales, oligopolio, Adif, barreras: infraestructura, inversión,
    regulación)
  \item
    2.2 Evolución de la demanda:(pasajeros, crecimiento tras
    liberalización, datos INE, gráfico en R)
  \item
    2.3 Grado de concentración del mercado (cuotas de mercado, índice
    HHI, cálculo en R, interpretación)
  \item
    2.4 Comportamiento estratégico de las empresas (competencia en
    precios - Bertrand)
  \item
    2.5 Poder de mercado - (precios - ¿Renfe como empresa privilegiada
    dentro del oligopolio?)
  \end{itemize}
\item
  \begin{enumerate}
  \def\labelenumi{\arabic{enumi}.}
  \setcounter{enumi}{2}
  \tightlist
  \item
    DISCUSIÓN Y CONCLUSIONES
  \end{enumerate}
\item
  \begin{enumerate}
  \def\labelenumi{\arabic{enumi}.}
  \setcounter{enumi}{3}
  \tightlist
  \item
    REFERENCIAS BIBLIOGRÁFICAS\\
    (INE, CNMC, informes sectoriales, manual de microeconomía)
  \end{enumerate}
\end{itemize}

\subsection{Otras cuestiones}\label{otras-cuestiones}

¿Cómo cito en el texto y cuando? Preguntar por: Leyes, história\ldots{}

\subsection{Fuentes}\label{fuentes}

\emph{falta citar e incluir las referencias en el texto}

\begin{itemize}
\tightlist
\item
  https://www.transportes.gob.es/el-ministerio/campanas-de-publicidad/2021-anio-europeo-del-ferrocarril/conociendo-el-ferrocarril/12-hitos
\item
  https://www.boe.es/buscar/doc.php?id=DOUE-L-1991-81196
\item
  https://www.canalsur.es/noticias/andalucia/la-red-de-alta-velocidad-espanola-la-mas-extensa-de-europa-con-4000-kilometros/2236223.html?utm\_source=chatgpt.com
\end{itemize}




\end{document}
